\section*{Заключение}
\addcontentsline{toc}{section}{Заключение}

В ходе выполнения курсового проекта было разработано веб-приложение для централизованного управления и мониторинга CI/CD-пайплайнов, функционирующее по принципу единого информационного окна («Single Pane of Glass»). Разработанная система решает проблему фрагментации данных при мультирепозиторной разработке и использовании гетерогенных систем автоматизации, объединяя сбор и отображение метрик из платформ GitHub Actions и GitLab CI/CD.

В процессе выполнения работы были решены следующие задачи:
\begin{itemize}
    \item проведен анализ предметной области и существующих решений автоматизации (GitHub Actions, GitLab CI/CD, Jenkins, Argo CD), сформулированы критерии сравнения и выявлена потребность в специализированном легковесном диспетчере сборок;
    \item обоснован и выбран технологический стек, включающий библиотеку React и TypeScript для клиентской части, асинхронный фреймворк FastAPI на языке Python для серверной части и реляционную СУБД PostgreSQL для долговременного хранения данных;
    \item составлены функциональные и нефункциональные требования к системе, разработаны диаграмма вариантов использования (Use Case) и пользовательские истории (User Stories);
    \item спроектирована реляционная модель данных и физическая схема базы данных в третьей нормальной форме, включающая механизмы индексации и безопасного хранения зашифрованных токенов доступа;
    \item реализована серверная архитектура веб-приложения с модулем аутентификации на основе JWT, асинхронными HTTP-адаптерами к внешним REST API и алгоритмом симметричного шифрования Fernet;
    \item разработан интерактивный SPA-интерфейс дашборда на React, обеспечивающий фильтрацию сборок, отображение карточек статусов конвейеров, просмотр консольных журналов ошибок и повторный запуск упавших задач в один клик;
    \item проведено функциональное тестирование основных пользовательских сценариев и выполнено развертывание сервиса.
\end{itemize}

Все поставленные в курсовом проекте цели и задачи успешно выполнены. Разработанный программный комплекс демонстрирует стабильную работу, сокращает время реакции инженеров на сбои в пайплайнах и может быть расширен за счет добавления новых адаптеров CI/CD-платформ.

Исходный код проекта опубликован и доступен в репозитории:
\begin{itemize}
    \item репозиторий проекта: \url{https://github.com/h3sug0/coursework-web-app-ci-cd};
    \item веб-приложение, развернутое на хостинге: \url{https://coursework-web-app-ci-cd.onrender.com}.
\end{itemize}